% Generated draft: requires mathematical and source review.
\documentclass[11pt]{article}
\usepackage[margin=25mm]{geometry}
\usepackage{fontspec}
\setmainfont{Latin Modern Roman}
\usepackage{amsmath,amssymb,mathtools,unicode-math}
\setmathfont{Latin Modern Math}
\usepackage{enumitem,needspace,xurl,hyperref}
\setlength{\parindent}{0pt}
\setlength{\parskip}{5pt}
\setlength{\emergencystretch}{3em}
\newcommand{\N}{\mathbb{N}}
\newcommand{\Z}{\mathbb{Z}}
\newcommand{\Q}{\mathbb{Q}}
\newcommand{\R}{\mathbb{R}}
\newcommand{\C}{\mathbb{C}}
\newcommand{\F}{\mathbb{F}}
\newcommand{\A}{\mathbb{A}}
\newcommand{\PP}{\mathbb{P}}
\newcommand{\E}{\mathbb{E}}
\newcommand{\eps}{\varepsilon}
\newcommand{\dd}{\,\mathrm{d}}
\begin{document}
\section*{UnsolvedMath ID: 6700083}
\textbf{Source problem code:} catalogue entry
\subsection{Definitions and mathematical statement}
For a smooth \(n\)-manifold \(X\), let \(\mathcal R^0(X)\) be its continuous positive definite Riemannian metrics. The source uses a volumic scalar curvature condition \(\operatorname{Sc}^{\mathrm{vol}}_n(g)>0\), whose definition is not included. The proposed density statement is that
\[
 \{g\in\mathcal R^0(X):\operatorname{Sc}^{\mathrm{vol}}_n(g)>0\}
\]
is dense, in the intended \(C^0\) topology, among all Riemannian metrics on every \(n\)-manifold \(X\), for \(n\geq3\). Thus each admissible target metric \(g_0\) should be a \(C^0\) limit of continuous metrics \(g_j\) satisfying that volumic positivity condition. The text's footnote 69 explicitly cautions that ``it may be safer to assume \(n\geq4\)'' and says conformal representation of surface metrics makes the analogous approximation unlikely in dimension two. This caveat is part of the proposal; dimension three is not silently removed or certified.

For comparison, the source reports Lohkamp's theorem [82]: \(C^2\) metrics with \(\operatorname{Scal}\leq-1\) are \(C^0\)-dense among Riemannian metrics in dimensions \(n\geq3\). This comparison is a reported theorem, not a replacement of volumic positivity by ordinary positive scalar curvature.

The passage also tentatively proposes extending many properties of smooth positive scalar curvature manifolds to spaces with benign singularities, for example Alexandrov spaces with sectional curvature bounded below by \(-1\). Its stated comparison convention is: a geodesic metric space \(X\) has that lower bound if each quadruple \(x_1,x_2,x_3,x_4\) admits points \(x'_1,x'_2,x'_3,x'_4\) in the hyperbolic plane of curvature \(-1\) such that
\[
 d_{\mathbb H^2}(x'_i,x'_j)=d_X(x_i,x_j)\ (1\leq i,j\leq3),\qquad
 d_{\mathbb H^2}(x'_i,x'_4)\leq d_X(x_i,x_4)\ (1\leq i\leq3).
\]
Geodesic means that any two points are joined by a path realizing their distance. No particular positive scalar curvature property or corresponding singular-space scalar curvature definition is specified, so this extension remains a qualitative proposal. The next expected result announced at the end is absent.
\subsection{Short English statement}
Are continuous metrics with volumically positive scalar curvature dense in the \(C^0\) topology in dimensions at least three, with an explicit source caution about dimension three? The passage also proposes extending positive scalar curvature theory to suitable Alexandrov spaces, without stating a particular extension theorem.
\subsection{Sources}
\begin{itemize}
\item Local UnsolvedMath registry, permanent ID 6700083, source code catalogue entry; Misha Gromov, 2017; updated 2026-08-06.
\item Gromov, \emph{101 Questions, Problems and Conjectures around Scalar Curvature} (2017), printed tag [?85], occurrence 86, PDF page 83, including footnotes 69 and 70: \url{https://www.ihes.fr/~gromov/wp-content/uploads/2018/08/101-problemsOct1-2017.pdf}. Registry access date 2026-07-29; normalized extraction from a public author-hosted incomplete PDF; current status and reuse terms need review.
\item Source reference [82] is attributed to Lohkamp for the negative scalar curvature density theorem, with no complete bibliographic entry supplied.
\item Source reference [4] is associated in the extraction with Alexander, Kapovitch, and Petrunin, \emph{Alexandrov geometry} (2017), and further references; complete bibliographic details are absent.
\item The supplied partial-progress summary reports supporting examples and broader flexibility results without exact further references or a verified full resolution.
\end{itemize}
\end{document}
